% Figure 1: Paradigm Shift Diagram
\begin{figure}[htbp]
\centering
\begin{tikzpicture}[
    scale=0.85,
    transform shape,
    box/.style={
        rectangle,
        rounded corners=3pt,
        draw=#1!70!black,
        fill=#1!20,
        thick,
        text width=5.5cm,
        minimum height=1.8cm,
        align=center,
        font=\small
    },
    arrow/.style={
        ->,
        >=stealth,
        thick,
        #1
    }
]

% Left Panel: Prediction Paradigm
\node[font=\large\bfseries, color=blue!70!black] at (-4, 6) {Prediction Paradigm};
\node[font=\footnotesize, color=gray] at (-4, 5.5) {(1980s-2020)};

\node[box=blue] (input1) at (-4, 4) {
    \textbf{Input: Known Proteins}\\
    Sequences/Structures
};

\node[box=blue!60] (analysis) at (-4, 1.5) {
    \textbf{Computational Analysis}\\
    Sequence similarity\\
    Structure comparison\\
    ML classification
};

\node[box=blue!40] (output1) at (-4, -1) {
    \textbf{Output: Classification}\\
    ``Will these interact?''\\
    Answer: YES / NO
};

\draw[arrow=blue!70] (input1) -- (analysis);
\draw[arrow=blue!70] (analysis) -- (output1);

\node[font=\footnotesize\itshape, color=blue!60] at (-4, -2.5) {Mode: Discovery \& Understanding};

% Right Panel: Design Paradigm
\node[font=\large\bfseries, color=green!70!black] at (4, 6) {Design Paradigm};
\node[font=\footnotesize, color=gray] at (4, 5.5) {(2021-Present)};

\node[box=green] (input2) at (4, 4) {
    \textbf{Input: Design Goal}\\
    ``What interaction do we WANT?''
};

\node[box=green!60] (design) at (4, 1.5) {
    \textbf{Generative AI Design}\\
    Diffusion models\\
    Inverse folding\\
    Multi-objective optimization
};

\node[box=green!40] (output2) at (4, -1) {
    \textbf{Output: Novel Proteins}\\
    ``Proteins that WILL bind''\\
    (Candidates: 100s-1000s)
};

\node[box=green!30] (validate) at (4, -3.5) {
    \textbf{Validation \& Iteration}\\
    Computational screening\\
    Experimental testing
};

\draw[arrow=green!70] (input2) -- (design);
\draw[arrow=green!70] (design) -- (output2);
\draw[arrow=green!70] (output2) -- (validate);
\draw[arrow=green!70, rounded corners=8pt] (validate.west) -- ++(-2,0) |- (input2.west);

\node[font=\footnotesize\itshape, color=green!60] at (4, -5) {Mode: Creation \& Engineering};

% Center Arrow
\draw[->, ultra thick, red!70!black] (-0.3, 1.5) -- (0.3, 1.5);
\node[font=\small\bfseries, color=red!70!black, rotate=90] at (0, 1.5) {SHIFT};

\end{tikzpicture}
\caption{The paradigm shift from prediction to design in AI-driven protein-protein interaction research. Left: The prediction paradigm (1980s-2020) focuses on analyzing existing proteins to predict interactions. Right: The design paradigm (2021-present) enables creation of novel proteins with desired interaction properties through iterative design-test cycles.}
\label{fig:paradigm_shift}
\end{figure}
